\documentclass[10pt]{article}
\usepackage[utf8]{inputenc}
\usepackage[T1]{fontenc}
\usepackage{graphicx}
\usepackage[export]{adjustbox}
\graphicspath{ {./images/} }
\usepackage{amsmath}
\usepackage{amsfonts}
\usepackage{amssymb}
\usepackage[version=4]{mhchem}
\usepackage{stmaryrd}
\usepackage{bbold}

\DeclareUnicodeCharacter{0131}{$\imath$}

\begin{document}
отклоняющимся аргументом, дифференциальные уравнения с частными производными). Имеется и ряд результатов, характеризующих соответствующие свойства. Для нелинейных систем известен лишь ряд локальных теорем о наблюдаемости. Решения проблемы наблюдения нашли многочисленные применения в задачах синтеза в условиях неполной информации о координатах, в том числе в задачах оптимальной стабилизации (см. [3]-[5], [14], [15]).

Задача синтеза становится особенно содержательной, когда информация об уравнениях управляемого процесса, исходных начальных условиях и текущих параметрах искажена возмущениями. Если описание возмущений носит статистич. характер, то задачи оптимального управления рассматривают в рамках тео-

\includegraphics[max width=\textwidth, center]{2023_07_08_a64f6fd79774dd23de19g-1}
у п р а в е н и я. Әта теория, начатая с решения стохастических программных задач [16], в наибольшей степени разработана для систем вида

\[
\dot{x}=f(t, x, u)+g(t, x, u) \eta, x\left(t_{0}\right)=x^{0},
\]

со случайными возмущениями $\eta(t)$, описываемыми гауссовскими диффузионными процессами или более общими классами марковских процессов (начальный вектор обычно также считают случайным). При этом, как правило, предполагается, что заданы нек-рые вероятностные характеристики величин $\eta$ (напр., сведения о моментах соответствующих распределений или о параметрах стохастич. уравнений, описывающих әволюцию процесса $\eta(t)$ ).

В общем случае применение программных и синтезирующих управлений здесь дает существенно различные значения оптимальных показателей $J$ качества (роль таких показателей могут играть, напр., те или иные средние оценки неотрицательных функционалов, заданных на траекториях процесса). Задача синтеза стохастического оптимального управления теперь имеет очевидные преимущества, т. к. непрерывное измерение координат системы нозволяет корректировать движение с учетом реального хода случайного процесса, непредсказуемого заранее. Здесь было обнаружено, что метод динамич. программирования совершенно естественным образом сопрягается с теориеӥ бесконечно малых производящих операторов для полугрупп преобразований, генерируемых марковскими случайными процессами. Эти обстоятельства позволили построить и строго обосновать серию достаточных условий оптимальности, приведших к решению на конечном и бесконечном интервале времени ряда задач о синтезе стохастического оптимального управления $\mathrm{c}$ полной и неполной информацией о текущих координатах, стохастич. задач преследования и т. Д. При этом существенно, что для справедливости принципа оптимальности управление $u$ здесь должно строиться в каждый момент времени $t$ как функция "достаточных координат» $z$ процесса, для к-рых будет обеспечено свойство марковости (см. [5], [6], [17], [18]).

Таким путем, в частности, была разработана теория $о$ п т ІІ м а ль н о й с т ох а с т и ч е с к о й с т а б пі л и з а ц и и, связанная с соответствующей теориеӥ устоӥчивости по Ляпунову, развитой для стохастич. систем [19].

Для формирования синтезıрующего оптимального управления, а также для иных целей управления целесообразно оценивать состояние стохастич. системы по результатам измерения. Решению этого вопроса при условии, что процесс измерения искажается вероятностными «шумами» (т. е. решению задачи наблюдения в условия $\mathrm{x}$ случайных возмущений), посвящена те ор ия с тохастическ о ї би ль т р ации. Наиболее полные репения здесь известны для линейных систем с квадратичными критериями оптимума (т. н. фи л ь т р К а л м а на-Б ь го с п, см. [13]). В применении әтої теории к задаче синтеза стохастического оптимального управления были выявлены условия, обеспечивающие справедливость принципа разделения, позволяющего решать собственно задачу управления независимо от задачи оценивания текущих позиций на основе достаточных координат процесса (см. [20]; более общим закономерностям стохастич. фильтрации, а также задачам стохастич. оптимального управления, когда само управлевие выбирается в классө марковских процессов диффузионного типа, посвящены работы [18], [21]).

Строго формализованное решение задачи стохастического оптимального управления неизбежно соприкасается с проблемой корректного обоснования вопросов существования решений соответствующих стохастических дифференциальных уравнений. Последнее обстоятельство порождает определенные трудности в решении задач стохастического оптимального управления при валичии неклассич. ограничений.

Содержательный процесс динамич. оптимизации вовникает в задачах синтеза оптимального управления в условиях неопределенности (см. Оптимальное управление програмяное). Позиционные решения в общем случае позволяют и здесь улучшить показатели качества процесса по сравнению с программными, представляющими собой все-таки результат статич. оптимизации (проводимоӥ, правда, в пространстве динамич. систем и функций-управлений). К решению задач тогда привлекают понятия и методы теории игр.

Пусть имеется система

\[
\dot{x}=f(t, x, u, w), t_{0} \leqslant t \leqslant t_{1},
\]

при ограничения

\[
x^{0}=x\left(t_{0}\right) \in X^{0} \subseteq \mathbb{R}^{n}, u \in U \subseteq \mathbb{R}^{p}, w \in W \subseteq \mathbb{R}^{q},
\]

на начальный вектор $x^{0}$, управление $u$ и возмущения $w$. В отличие от игрока-союзника, представляющего исходное управление $u$, подлежащее определению, воздействия $w$ здесь трактуют как управление игрокапротивника и допускают к рассмотрению любые целенаправленные стратегии $w$, формируемые на основе любой допустимой информации. При әтом цели управления могут быть сформулированы с точки зрения каждого из игроков в отдельности. Если названные цели противоположны, то возникает задача к о нФл и к т н о г о у п р а в ле н и я. Исследование задач позиционного управления в условия х конфликта или неопределенности составляет предмет теории $\partial и \phi-$ беренциальных иәр.

Процесс формирования позиционного оптимального управления в условиях неопределенности может быть осложнен неполнотой информапии о текущем состоявии. Так, в системе (8) могут быть доступны лить результаты косвенных измерений фазового вектора $x$ реализации $y[t]$ функции

\[
y(t)=g(t, x, \xi)
\]

где неопределенные параметры $\xi$ стеснены известным априорным ограничением $\xi \in \Xi$. Знание $y[t], t_{0} \leqslant t \leqslant \vartheta$ (при заданном $u(t)$ ), позволяет построить в фазовом пространстве и нфо р м а и он н ую о б л а с т $X(\vartheta, y(\cdot))$ состояний системы (8), совместимых с реализацией $y[t]$, уравнением (9) и ограничениями на $w$, $\xi$. Среди әлементов $X(\vartheta, y(\cdot))=X(\grave{\vartheta}, \cdot)$ будет содержаться и неизвестное истинное состояние системы (8), к-рое может быть оценено выбором век-рой точки $x^{*}(\vartheta, \cdot)$ из $X(\vartheta, \cdot)$ (напр., «центра тяжести» или «чебышевского центра» $X(\vartheta, \cdot))$. Изучение әволюции областеӥ $X(\vartheta, \cdot)$ и динамики векторов $x^{*}(\vartheta, \cdot)$ составляет содержавие

\begin{center}
\includegraphics[max width=\textwidth]{2023_07_08_a64f6fd79774dd23de19g-1(1)}
\end{center}


\end{document}